% =============================================================================
% Project : THE_PHILOSOPHY_OF_VIABILITY
% Book    : Book I — Ontology of Viability
% Book ID : VS.B01
% Part    : Part I — Identity and Manifestation (VS.B01.P01)
% Chapter : Chapter 02 — Invariance
% Chapter ID : VS.B01.P01.C02
%
% File    : 4_ASSETS/diagrams/VS.B01.P01.C02.S10.SS11.P004.DG01.scope_of_an_invariance_claim.tex
%
% Section / Chapter / Asset:
%   Diagram DG01 — Scope of an Invariance Claim
%
% Structural Identifier : VS.B01.P01.C02.S10.SS11.P004.DG01
% Asset Type : DG
%
% Purpose :
%   To display the elements that delimit an invariance claim and distinguish
%   its warranted conclusion from conclusions that require independent
%   ontological justification.
%
% Representational limit :
%   A justified preservation finding is restricted to its specified subject,
%   determination, transformation, domain, respect, criterion, and
%   conditions. Further ontological conclusions require independent
%   justification.
%
% Dependencies :
%   TikZ; libraries used: arrows.meta, calc, positioning, fit. All are loaded
%   by the project preamble; no per-file package or library loading.
%
% =============================================================================

\begin{figure}[htbp]
    \centering
    \hyphenpenalty=10000 \exhyphenpenalty=10000\relax  % no hyphenation inside diagram labels

    \begin{tikzpicture}[
        font=\small,
        node distance=0.35cm,
        % --- base box ---
        box/.style={
            draw, rounded corners=3pt, align=center,
            inner sep=5pt, line width=0.7pt
        },
        % --- node roles ---
        specbox/.style={box, text width=3.1cm, minimum height=1.1cm,
                        fill=blue!8, draw=blue!55!black},
        assessbox/.style={box, text width=5.6cm, minimum height=1.0cm,
                          fill=blue!16, draw=blue!55!black},
        resultbox/.style={box, text width=5.6cm, minimum height=1.1cm,
                          font=\small\bfseries, line width=0.9pt,
                          fill=green!8, draw=green!40!black},
        limitbox/.style={box, text width=3.1cm, minimum height=1.1cm,
                         dashed, fill=gray!10, draw=black!60},
        % --- frames and labels ---
        vsgroup/.style={draw=black!45, dashed, rounded corners=5pt,
                        inner sep=8pt, line width=0.5pt},
        grouplabel/.style={font=\footnotesize\itshape, text=black!75,
                           align=center},
        % --- arrows and lines ---
        directed/.style={-{Stealth[length=2.5mm]}, line width=0.9pt},
        vslink/.style={draw=black!55, dashed, line width=0.7pt},
        boundary/.style={draw=black!70, line width=0.9pt,
                         dash pattern=on 6pt off 2pt on 1pt off 2pt}
    ]

    % ---------------------------------------------------------------
    % Claim specification: six elements, no order among them
    % ---------------------------------------------------------------
    \node[specbox] (subject)        {Existent / Subject};
    \node[specbox, right=of subject]        (determination)  {Specified Determination};
    \node[specbox, right=of determination]  (transformation) {Transformation};
    \node[specbox, below=of subject]        (domain)         {Domain and Respect};
    \node[specbox, below=of determination]  (criterion)      {Preservation Criterion};
    \node[specbox, below=of transformation] (conditions)     {Conditions and Limits};

    \node[vsgroup, fit=(subject)(determination)(transformation)
                       (domain)(criterion)(conditions)] (spec) {};
    \node[grouplabel, above=0.1cm of spec] {Elements delimiting the claim};

    % ---------------------------------------------------------------
    % Assessment and warranted result
    % ---------------------------------------------------------------
    \node[assessbox, below=0.8cm of spec] (assessment) {
        \textbf{Specified Invariance Assessment}
    };
    \node[resultbox, below=0.8cm of assessment] (result) {
        Preservation Established Within Scope
    };

    \draw[directed] (spec)       -- (assessment);
    \draw[directed] (assessment) -- (result);

    % ---------------------------------------------------------------
    % Inferential boundary
    % ---------------------------------------------------------------
    \coordinate (bnd) at ($(result.south)+(0,-1.0)$);
    \draw[boundary] (spec.west |- bnd) -- (spec.east |- bnd);
    \node[grouplabel, fill=white, inner sep=2pt] at ($(bnd)+(-1.75,0)$) {
        Inferential Boundary
    };

    % ---------------------------------------------------------------
    % Further ontological questions (same columns as the top grid)
    % ---------------------------------------------------------------
    \node[limitbox, below=1.0cm of bnd] (persistence) {
        Persistence and Continuity
    };
    \node[limitbox, left=of persistence]  (identity) {Numerical Identity};
    \node[limitbox, right=of persistence] (totality) {
        Complete Ontological Constitution
    };

    % Undirected connectors: not entailed by the invariance result
    \draw[vslink] (result.south) -- (persistence.north);
    \draw[vslink] ($(result.south)+(0,-0.4)$) -| (identity.north);
    \draw[vslink] ($(result.south)+(0,-0.4)$) -| (totality.north);

    % ---------------------------------------------------------------
    % Note
    % ---------------------------------------------------------------
    \node[below=0.7cm of persistence, text width=10cm, align=center,
          font=\footnotesize\itshape] {
        These further conclusions are not entailed by the
        invariance finding alone. Each requires independent
        ontological justification.
    };

    \end{tikzpicture}

    \caption{The scope of an invariance claim. A justified preservation
    finding is restricted to its specified subject, determination,
    transformation, domain, respect, criterion, and conditions.
    Further ontological conclusions require independent
    justification.}
    \label{fig:scope-of-an-invariance-claim}
\end{figure}

% =============================================================================
% END OF DIAGRAM DG01
% =============================================================================
